\chapter{Metodología técnica (CI/CD)}\label{cap:metodologia_tecnica}

Mi idea es haber explicado todas las herramientas en el capítulo 4, pero aquí se podría entrar más en detalle sobre cómo se han configurado y utilizado en lo que respecta a ci/cd, gestión de ramas, commits, etc. También se podrían incluir políticas de código, como pre-commits, hooks, plantillas de issues, etc. que no se hayan mencionado en capítulos anteriores.

\section{Políticas del repositorio}
Ramas, commits y versiones. Se puede mencionar los pre-commits, hooks, Github, Commitizen, Convenciones de commits, plantillas de issues, ... y su importancia en el desarrollo para trabajar de manera más organizada

\section{Integración y Despliegue Continuo (CI/CD)}
Descripción del ciclo ci/cd
\subsection{Configuración del Pipeline y automatización}
Implementación de flujos de trabajo en GitHub Actions y qué herramientas incluyen

\section{Perfiles de configuración y entornos}
Variables de entorno, secretos y configuración específica para los entornos de desarrollo, pruebas y producción.

\section{Métricas de calidad y análisis estático de código}
SonarCloud para el análisis estático de código, deuda técnica y cobertura de pruebas. (Para ello también habría que ver hasta qué punto se ha hablado de Sonar en capítulos anteriores)

Aqui quiero hacer especial énfasis para explicar la importancia de estas métricas, cómo se han configurado y qué resultados han dado, mostrando gráficas o ejemplos concretos de cómo han ayudado a mejorar la calidad del código y a identificar áreas de mejora.
También es importante recalcar que su uso es vital para evitar deuda técnica, y en el capítulo 9 de resultados diré que me arrepiento de haberlo implementado tan tarde y que me ha hecho tener que aplazar un par de issues que tenias planteadas para la entrega del tfg. (también podría ser en el capítulo de conclusiones)

